
{\fontsize{10.25}{12.1}\selectfont
La maladie épidémique qui règne depuis environ cinq mois, parmi les gens de la campagne à Meyrueis et ses environs, Lanuéjols, Montjardin, Serviguières et quelques autres lieux moins considérables, est une fièvre, maligne, bilieuse, qui, à ce qu’on affirme, a fait de grands ravages à Séverac en Rouergue et a parcouru le Vallon du Tarn. Cette maladie prélude pendant plusieurs jours des malaises, des lassitudes, des pesanteurs d’estomac, le dégoût et le cours de ventre bilieux... la langue auparavant rouge, sèche, se racornit, se fend en différents sens et se retire dans le fond de la bouche, ou se couvre d’une croûte noirâtre et raboteuse qui s’étend sur les gencives et les lèvres qu’elle fait paraître brûlées et dans cet état ils ne peuvent ni parler ni avaler à moins qu’on ait l’attention d’humecter souvent ces parties.

Il est très douteux qu’on ait vus des bubons.

Le terme le plus ordinaire est du quatorzième ou quinzième jour, parfois 22 ou 23.

.

Les enfants en ont tous réchappé mais à Lanuéjols il y a eu 45 morts.

\medskip
\textit{Quelles sont les causes les plus apparentes de cette épidémie ?}

\medskip
« L’excessive misère qui règne depuis longtemps parmi les habitants de ces montagnes. La mauvaise nourriture et la malpropreté qui en sont inséparables, les alternatives de froid et de chaud, les orages mais surtout les brouillards qu’ils y ont presque continuellement essuyés. Une longue suite de mauvaises récoltes a rendu cette misère si sensible qu’à ne considérer que l’état actuel de ces infortunés, on ne saurait se persuader qu’ils aient jamais pu être les pères nourriciers de la plus grande partie des Cévennes...

\medskip
... relégués dans le fond des vallons, tous couverts de haillons et ayant à peine une chemise de rechange, ils y habitent de misérables chaumières placées sur le bord de quelque torrent et entourées de fumiers de buis qui exhalent une odeur des plus insupportable....

\medskip
... obligé de pénétrer dans ces chaumières on est tout étonné de voir que l’air intérieur de ces réduits obscurs et malsains y est à grand peine renouvelé par l’ouverture étroite d’une ample cheminée, d’une porte écrasée et d’une fente qui tient lieu de fenêtre pratiquée sur l’un des côtés de cette porte et de n’y trouver que quelques mauvaises planches mal assurées dans le mur, dont ils n’ont maintenant que faire, mais qui, autrefois, leur servaient à ranger leurs denrées.

\medskip
... on y distingue également un vieux coffre où ils tiennent leur peu de provision, et une vaste caisse carrée remplie de feuilles de frêne couverte de quelques lambeaux de grosse toile et d’une mauvaise couverture, qui leur sert lieu de grabat sur lequel le vieillard décrépit et l’enfant nouvellement né, la fille et le garçon, le sain et le malade, le mourant et le mort sont souvent confondus....

\medskip
Leur nourriture ordinaire est de l’orge la plus grossière dont ils font un peu de pain, ou qu’ils font cuire avec le lard le plus rance qui leur sert ensuite de pitance.. On les a vus, en dernier lieu, faire le bouillon ou la soupe pour leurs malades avec le vieux oing*.

Les autres aliments dont ils usent sont les aulx, les oignons et, en hiver, quelques pommes de terre, des raves, les plus mauvais légumes et des fromages encore plus mauvais.

Aussi en est-il mort un très grand nombre.

\medskip
\textit{\textbf{Oing :} Vieille graisse de porc fondue servant à graisser un mécanisme.}
}

